\chapter{Introduction}\label{ch:introduction}

\section{Context}

Large-scale software-intensive systems maintained by industrial research
organisations are typically mission-critical, long-lived and subject to
continuous evolution \cite{systa:reverse}. Many operate in
safety-critical or performance-sensitive domains such as medical imaging
and industrial automation, where an incorrect change or a poorly
informed architectural decision carries real technical and economic
consequences. An accurate understanding of system behaviour is therefore
needed before any refactoring, modernisation or architectural
restructuring.

This understanding is captured by the \emph{Interaction Model}: a
formal, scenario-specific description of the message exchanges between
components, most commonly expressed as sequence diagrams in \ac{UML} or,
more recently, \sysml. In practice, Interaction Models rarely stay
current with the source code they describe. Complexity accumulates
from evolving requirements, heterogeneous technologies and concurrency
\cite{systa:reverse}, and diagrams accurate when drawn drift out of sync
as the system changes. Engineers are consequently left to reconstruct
system behaviour directly from source code and limited runtime
observations, a process that is time-consuming, error-prone and heavily
dependent on individual expertise \cite{systa:reverse}.

The problem is worse in event-driven systems, where asynchronous
message passing, callbacks and indirect control flow obscure how
components relate to one another \cite{cornelissen:systematic}.
Determining which components interact, in what order and under what
conditions, is non-trivial, yet essential for identifying architectural
bottlenecks, isolating faulty behaviour and assessing the impact of
proposed changes, which is why industry needs accurate,
scenario-specific interaction models reflecting how a system actually
behaves under concrete execution scenarios \cite{egyed:scenario}.

A concrete illustration of this problem surfaced during this project: a
real, previously-resolved defect in a large industrial \Cpp\ system
(\cref{sec:scenario-exercise}). A stale proxy object was never freed
because a reset call was missing from the connection-handling code, and
the system silently stopped delivering change-notification events after
a service restart.

The \ac{GDI} project at \ac{TNO}-\ac{ESI} addresses this problem using
graph-based program analysis to support automatic extraction and
maintenance of software architecture models. This thesis contributes a
hybrid approach combining two complementary sources of evidence: a
static call graph extracted from source code, and timestamped execution
traces captured at runtime.

\section{Problem Statement}

Static analysis reveals call dependencies and structural relationships,
giving a complete view of potential interactions, but it typically
over-approximates behaviour and cannot distinguish interactions specific
to one scenario from those occurring in every other scenario
\cite{systa:reverse}. Trace-based analysis is scenario-specific by
construction, but inherently incomplete \cite{cornelissen:systematic}:
in industrial systems, traces are often sparse or noisy due to limited
logging, instrumentation overhead, or execution paths simply not
exercised during recording.

Hybrid static-dynamic approaches have been proposed
\cite{labiche:hybrid, cornelissen:systematic}, but existing methods
primarily use static structure to organise or contextualise an
already-observed trace (\cref{ch:related}); they do not actively infer
interactions the trace never captured. No existing approach
systematically combines the two directions at once: using static
structure to recover what a trace missed, and the trace to trim an
over-approximate static graph down to what is relevant to the scenario
at hand. Closing that gap in both directions is the core problem this
thesis addresses.

A further complication comes from the publish-subscribe indirection
common in event-driven \Cpp\ systems. \emph{CPSCore}, the open-source
framework used as the case study, routes inter-component communication
through an \ac{IDC} abstraction layer backed by \texttt{boost::signals2}
callbacks and Redis channels. A static call graph built from function
call sites alone contains edges \emph{into} the IDC layer, but cannot
reveal which component receives each published packet, since the
subscription topology is implicit in callback registrations rather than
explicit in call sites. Recovering that missing structure is part of
this thesis's contribution (\cref{sec:macro-resolution}).

\section{Research Questions}

This thesis is guided by three research questions:

\begin{description}
  \item[\rqone] Does combining a static call graph with runtime execution
    traces reconstruct scenario-specific Interaction Models more
    accurately than either source alone?

  \item[\rqtwo] Does guiding trace collection with the static call graph
    remove reconstruction noise without losing scenario coverage?

  \item[\rqthree] How well does a reconstructed \sysml Interaction Model
    consolidate static and dynamic evidence into a single artefact
    sufficient for system-behaviour comprehension?
\end{description}

Each is paired with a testable hypothesis, answered by a dedicated
experiment (\cref{ch:experiments,ch:analysis}):

\begin{description}
  \item[H1 (\rqone).] Combining static analysis and dynamic instrumentation
    improves reconstruction accuracy over either source alone.
  \item[H2 (\rqtwo).] Code-graph-guided instrumentation improves scenario
    focus by suppressing irrelevant execution detail relative to full,
    unguided instrumentation.
  \item[H3 (\rqthree).] Reconstructed interaction models improve agent task
    performance on scenario comprehension, relative to reasoning from raw
    source code alone.
\end{description}

\rqone, \rqtwo\ and \rqthree\ are answered empirically by the H1, H2 and
H3 experiments on the CPSCore case study (\cref{ch:experiments,ch:analysis}).

\section{Approach}

The approach consists of six stages, implemented as a chain of
composable agent skills. First, \textbf{static extraction}: a call graph
is extracted from the target system's \Cpp\ source and stored in a \neo
\acl{GDB} using the schema in \cref{sec:pg-schema}; for pub-sub systems
routing messages through callback registration, such as CPSCore's IDC
layer, this stage additionally resolves the implicit routing topology
(\cref{sec:macro-resolution}). Second, \textbf{code graph query}: a
scenario description is translated into a Cypher query and executed
against the static property graph, retaining only the interactions whose
caller and callee belong to the scenario's declared components. Third,
\textbf{dependency matrix extraction}: the scenario-scoped subgraph is
flattened into an ordered table of static call edges. Fourth,
\textbf{trace instrumentation and runtime tracing}: probes are injected
into the source by the instrumentation tool (\cref{sec:csp-matcher}), the
instrumented program is run, and the emitted timestamped trace lines are
normalised into the same table format; this dynamic pipeline is kept
separate from the static property graph, and the two sources are
combined only when the reconstructed models are compared during
evaluation. Fifth, \textbf{agentic synthesis}: an \ac{LLM} is given the
scenario-scoped static and dynamic evidence as two separately labelled,
un-merged sources and weighs which single-source signals are genuine
against which are noise to decide the final set of architecturally
significant interactions (\cref{sec:agentic-synthesis}), the only stage
in the pipeline that is not purely mechanical. Sixth, \textbf{SysML v2
generation}: the synthesised
interaction set is projected into \sysml sequence diagrams by a
template-driven generator and written to a \texttt{.sysml} file.

The approach is developed and validated on \emph{CPSCore}, an
open-source event-driven \Cpp\ framework, across four scenarios spanning
happy-path to a constructed indirection-boundary case
(\cref{ch:experiments,ch:analysis}).

\section{Outline}

\Cref{ch:background} introduces the technical concepts underpinning the
approach: \sysml Interaction Models, property graphs, static call-graph
analysis, execution tracing and scenario-scoped querying.
\Cref{ch:related} surveys existing work on program comprehension,
architecture recovery and trace-based diagram synthesis. \Cref{ch:methods}
details the design: the property graph schema, the trace-collection
pipeline and the scenario-scoped querying procedure. \Cref{ch:experiments}
presents the implementation and results on the CPSCore case study.
\Cref{ch:analysis} analyses those results against the research questions
and hypotheses, and discusses threats to validity. \Cref{ch:conclusion}
summarises the contributions, discusses future work and concludes the
thesis.
